% GENERATED FILE. Edit canonical lessons, then run npm run generate:books.
% canonical-source-hash: fnv1a64:2557ec668a83121f
% canonical-lessons: SA-C418-jalam, SA-C418-badisam, SA-C418-suktih, SA-C418-mauktikam, SA-C418-pravalam

\chapter{Nets, Shells and Pearls}
\label{ch:sa-nets-shells-and-pearls}
\hlchaptermodality{418}

\begin{chapteropening}
\textbf{By the end of this chapter:} \emph{I can talk about what fishermen and divers bring up.}

Complete the last lesson of chapter 418: I can talk about what fishermen and divers bring up.
\end{chapteropening}

\section[\texorpdfstring{\sk{जालम्} (jālam)}{जालम् (jālam)}]{\sk{जालम्} (jālam) — a net}

\label{lesson:SA-C418-jalam}



\begin{quote}
Before the new one: say the Sanskrit for dry land, then the Sanskrit for the coastline.
\end{quote}

\subsection*{You'll want to know: \sk{जालम्}}
\textbf{\sk{जालम्}} — \emph{jālam} — \textquotedblleft{}a net\textquotedblright{}.

\begin{quote}
\sk{धीवरः} \sk{नद्यां} \sk{जालं} \sk{क्षिपति।}
\end{quote}

\begin{quote}
\emph{dhīvaraḥ nadyāṃ jālaṃ kṣipati.}
\end{quote}

\begin{quote}
The fisherman casts a net into the river.
\end{quote}

A neuter noun.

\begin{cousinweb}
\textbf{\sk{जालम्}} (\emph{jālam}) is the word MW gives for a net for catching birds or fish; MW gives no derivation, so its origin is uncertain (MW).
\end{cousinweb}

\begin{grammarlens}[title={What you've built}]
One more word for getting around.
\end{grammarlens}

\subsection*{Your turn}
\begin{itemize}
\raggedright
  \item \emph{Say it:} \emph{jālam}
  \item \emph{Say it:} \emph{jālam}, once more
  \item \emph{Say it:} \emph{dhīvaraḥ nadyāṃ jālaṃ kṣipati.} — The fisherman casts a net into the river.
  \item \emph{Recall:} say the Sanskrit for dry land, then the Sanskrit for the coastline, then say \emph{jālam} again
\end{itemize}

\begin{tcolorbox}[breakable,colback=teal!4,colframe=teal!35!black,title={Before you move on}]
What does \sk{जालम्} mean? (\textquotedblleft{}A net\textquotedblright{}.) Say it once more.
\end{tcolorbox}
\section[\texorpdfstring{\sk{बडिशम्} (baḍiśam)}{बडिशम् (baḍiśam)}]{\sk{बडिशम्} (baḍiśam) — a fishhook}

\label{lesson:SA-C418-badisam}



\begin{quote}
Before the new one: say the Sanskrit for the coastline, then the Sanskrit for a net.
\end{quote}

\subsection*{You'll want to know: \sk{बडिशम्}}
\textbf{\sk{बडिशम्}} — \emph{baḍiśam} — \textquotedblleft{}a fishhook\textquotedblright{}.

\begin{quote}
\sk{बालः} \sk{बडिशेन} \sk{मत्स्यं} \sk{गृह्णाति।}
\end{quote}

\begin{quote}
\emph{bālaḥ baḍiśena matsyaṃ gṛhṇāti.}
\end{quote}

\begin{quote}
The boy catches a fish with a hook.
\end{quote}

A neuter noun.

\begin{cousinweb}
\textbf{\sk{बडिशम्}} (\emph{baḍiśam}) is the word MW gives for a hook, a fishhook, also written \emph{vaḍiśa}; MW gives no derivation, so its origin is uncertain (MW).
\end{cousinweb}

\begin{grammarlens}[title={What you've built}]
One more word for getting around.
\end{grammarlens}

\subsection*{Your turn}
\begin{itemize}
\raggedright
  \item \emph{Say it:} \emph{baḍiśam}
  \item \emph{Say it:} \emph{baḍiśam}, once more
  \item \emph{Say it:} \emph{bālaḥ baḍiśena matsyaṃ gṛhṇāti.} — The boy catches a fish with a hook.
  \item \emph{Recall:} say the Sanskrit for the coastline, then the Sanskrit for a net, then say \emph{baḍiśam} again
\end{itemize}

\begin{tcolorbox}[breakable,colback=teal!4,colframe=teal!35!black,title={Before you move on}]
What does \sk{बडिशम्} mean? (\textquotedblleft{}A fishhook\textquotedblright{}.) Say it once more.
\end{tcolorbox}
\section[\texorpdfstring{\sk{शुक्तिः} (śuktiḥ)}{शुक्तिः (śuktiḥ)}]{\sk{शुक्तिः} (śuktiḥ) — an oyster shell}

\label{lesson:SA-C418-suktih}



\begin{quote}
Before the new one: say the Sanskrit for a net, then the Sanskrit for a fishhook.
\end{quote}

\subsection*{You'll want to know: \sk{शुक्तिः}}
\textbf{\sk{शुक्तिः}} — \emph{śuktiḥ} — \textquotedblleft{}an oyster shell\textquotedblright{}.

\begin{quote}
\sk{समुद्रस्य} \sk{तीरे} \sk{एका} \sk{शुक्तिः} \sk{अस्ति।}
\end{quote}

\begin{quote}
\emph{samudrasya tīre ekā śuktiḥ asti.}
\end{quote}

\begin{quote}
There is a shell on the seashore.
\end{quote}

A feminine noun.

\begin{cousinweb}
\textbf{\sk{शुक्तिः}} (\emph{śuktiḥ}) probably comes from the root √śuc, \textquotedblleft{}to shine\textquotedblright{}, as MW suggests: MW gives \emph{śukti} as a pearl-oyster or an oyster shell (MW).
\end{cousinweb}

\begin{grammarlens}[title={What you've built}]
One more word for getting around.
\end{grammarlens}

\subsection*{Your turn}
\begin{itemize}
\raggedright
  \item \emph{Say it:} \emph{śuktiḥ}
  \item \emph{Say it:} \emph{śuktiḥ}, once more
  \item \emph{Say it:} \emph{samudrasya tīre ekā śuktiḥ asti.} — There is a shell on the seashore.
  \item \emph{Recall:} say the Sanskrit for a net, then the Sanskrit for a fishhook, then say \emph{śuktiḥ} again
\end{itemize}

\begin{tcolorbox}[breakable,colback=teal!4,colframe=teal!35!black,title={Before you move on}]
What does \sk{शुक्तिः} mean? (\textquotedblleft{}An oyster shell\textquotedblright{}.) Say it once more.
\end{tcolorbox}
\section[\texorpdfstring{\sk{मौक्तिकम्} (mauktikam)}{मौक्तिकम् (mauktikam)}]{\sk{मौक्तिकम्} (mauktikam) — a pearl}

\label{lesson:SA-C418-mauktikam}



\begin{quote}
Before the new one: say the Sanskrit for a fishhook, then the Sanskrit for an oyster shell.
\end{quote}

\subsection*{You'll want to know: \sk{मौक्तिकम्}}
\textbf{\sk{मौक्तिकम्}} — \emph{mauktikam} — \textquotedblleft{}a pearl\textquotedblright{}.

\begin{quote}
\sk{शुक्तौ} \sk{मौक्तिकं} \sk{जायते।}
\end{quote}

\begin{quote}
\emph{śuktau mauktikaṃ jāyate.}
\end{quote}

\begin{quote}
A pearl forms inside an oyster.
\end{quote}

A neuter noun.

\begin{cousinweb}
\textbf{\sk{मौक्तिकम्}} (\emph{mauktikam}) is built on \emph{muktā}, \textquotedblleft{}a pearl\textquotedblright{}, with the ending \emph{-ika}; MW gives \emph{mauktika} as a pearl (MW).
\end{cousinweb}

\begin{grammarlens}[title={What you've built}]
One more word for getting around.
\end{grammarlens}

\subsection*{Your turn}
\begin{itemize}
\raggedright
  \item \emph{Say it:} \emph{mauktikam}
  \item \emph{Say it:} \emph{mauktikam}, once more
  \item \emph{Say it:} \emph{śuktau mauktikaṃ jāyate.} — A pearl forms inside an oyster.
  \item \emph{Recall:} say the Sanskrit for a fishhook, then the Sanskrit for an oyster shell, then say \emph{mauktikam} again
\end{itemize}

\begin{tcolorbox}[breakable,colback=teal!4,colframe=teal!35!black,title={Before you move on}]
What does \sk{मौक्तिकम्} mean? (\textquotedblleft{}A pearl\textquotedblright{}.) Say it once more.
\end{tcolorbox}
\section[\texorpdfstring{\sk{प्रवालम्} (pravālam)}{प्रवालम् (pravālam)}]{\sk{प्रवालम्} (pravālam) — coral}

\label{lesson:SA-C418-pravalam}



\begin{quote}
Before the new one: say the Sanskrit for an oyster shell, then the Sanskrit for a pearl.
\end{quote}

\subsection*{You'll want to know: \sk{प्रवालम्}}
\textbf{\sk{प्रवालम्}} — \emph{pravālam} — \textquotedblleft{}coral\textquotedblright{}.

\begin{quote}
\sk{समुद्रे} \sk{रक्तं} \sk{प्रवालं} \sk{भवति।}
\end{quote}

\begin{quote}
\emph{samudre raktaṃ pravālaṃ bhavati.}
\end{quote}

\begin{quote}
Red coral grows in the sea.
\end{quote}

A neuter noun.

\begin{cousinweb}
\textbf{\sk{प्रवालम्}} (\emph{pravālam}) probably comes from the root √val, \textquotedblleft{}to turn\textquotedblright{}, as MW suggests: MW gives \emph{pravāla} as a young shoot and as coral (MW).
\end{cousinweb}

\begin{grammarlens}[title={What you've built}]
One more word for getting around.
\end{grammarlens}

\subsection*{Your turn}
\begin{itemize}
\raggedright
  \item \emph{Say it:} \emph{pravālam}
  \item \emph{Say it:} \emph{pravālam}, once more
  \item \emph{Say it:} \emph{samudre raktaṃ pravālaṃ bhavati.} — Red coral grows in the sea.
  \item \emph{Recall:} say the Sanskrit for an oyster shell, then the Sanskrit for a pearl, then say \emph{pravālam} again
\end{itemize}

\begin{tcolorbox}[breakable,colback=teal!4,colframe=teal!35!black,title={Before you move on}]
What does \sk{प्रवालम्} mean? (\textquotedblleft{}Coral\textquotedblright{}.) Say it once more.
\end{tcolorbox}
